\lecture{11}{20. April 2026}{The De Saint Venant Problem, Pt. II}

\subsection{Finding the Barycenter of a Cross-Section}
The barycenter $G$, also called the \textit{centroid} for a homogeneous cross-section, is the ``geometric center'' of the area. 

I.e. if we define a local reference frame $\left\{ O; \Vec{e}_1, \Vec{e}_2 \right\}$ in the plane of the cross-section
Then every point in the cross-section will have a position vector:
\begin{equation*}
  \Vec{x} = x_1 \Vec{e}_1 + x_2 \Vec{e}_2
.\end{equation*}
The barycenter should then be:
\begin{equation*}
  \Vec{x}_G = \frac{1}{A} \int_A \Vec{x} \, \mathrm{d}A
.\end{equation*}
Or in component form as:
\begin{align*}
  x_{G1} &= \frac{1}{A}\int_A x_1 \, \mathrm{d}A \\
  x_{G2} &= \frac{1}{A} \int_A x_2 \, \mathrm{d}A
.\end{align*}

\subsection{Simple Bending: Only One Bending Moment Acts}
We now consider the special case, where $M_2 \neq 0$, is the only applied moment.
I.e. the body is only being bent about the $e_2$-axis.
Here, shear forces are zero, torsional moments are zero and hence only normal stress remains.

So,
\begin{equation*}
  \Vec{\tau}_3 = \Vec{0}, \qquad \sigma_{33} \neq 0
.\end{equation*}
That means the traction on the cross-section is purely axial:
\begin{equation*}
  \Vec{T}_{\Vec{e}_3} = \begin{pmatrix}
    0\\
  0\\
  \sigma_{33}\\
  \end{pmatrix}
.\end{equation*}

For Saint-Venant bending, compatibility leads to the fact that axial stress must be linear in the cross-section coordinates:
\begin{equation*}
  \sigma_{33} = a + bx_1 + cx_2
.\end{equation*}
Physically, this means the beam bends in such a way that one side is stretched and the other side is compressed.
The axial strain varies linearly across the section, so the axial stress also varies linearly.

Because $M_1 = 0$, the stress cannot vary with $x_2$, so the stress depends only on $x_1$:
\begin{equation*}
  \sigma_{33} = \alpha x_1
.\end{equation*}
There is no constant term here because there is no axial force $N$.
If there were a constant term, integrating over the area would give a nonzero axial resultant force.

So for pure bending about $e_2$:
\begin{equation*}
  N = 0, \qquad M_1 = 0, \qquad M_2 \geq 0
.\end{equation*}
Therefore $\sigma_{33}$ must be linear in $x_1$ only.

From the previous section, the moment resultant about $e_2$ was
\begin{equation*}
  M_2 = - \int_A x_1 \sigma_{33} \, \mathrm{d}A
.\end{equation*}
Inserting $\sigma_{33} = \alpha x_1$ yields
\begin{equation*}
  M_2 = - \int_A x_1 \left( \alpha x_1 \right)\, \mathrm{d}A = - \alpha \int_A x_1^2 \, \mathrm{d}A
.\end{equation*}
The integral
\begin{equation*}
  I_2 = \int_A x_1^2 \, \mathrm{d}A
\end{equation*}
is the second moment of area about $e_2$.
So
\begin{equation*}
  M_2 = - \alpha I_2
.\end{equation*}
Therefore:
\begin{equation*}
  \alpha = - \frac{M_2}{I_2}
.\end{equation*}
Substituting back into the stress expression gives:
\begin{equation*}
  \sigma_{33}(x_1) = - \frac{M_2}{I_2}x_1
.\end{equation*}

\subsubsection{General Case by Superposition}
Because the problem is linear, we can combine the simpler solutions.
If we have the axial force $N$ and the bending moments $M_1$ and $M_2$, then the axial stress is the sum of these three contributions.

A pure axial force gives uniform stress:
\begin{equation*}
  \sigma_{33}^{(N)} = \frac{N}{A}
.\end{equation*}
Here, every fibre carries the same normal stress.
From above, bending about $e_2$ results in
\begin{equation*}
  \sigma_{33}^{(M_2)} = - \frac{M_2}{I_2}x_1
.\end{equation*}
And similarly, bending about $e_1$ gives:
\begin{equation*}
  \sigma_{33}^{(M_1)} = \frac{M_1}{I_1}x_2
.\end{equation*}

So the total stress is:
\begin{equation*}
  \sigma_{33}(x_1, x_2) = \frac{N}{A} - \frac{M_2}{I_2}x_1 + \frac{M_1}{I_1}x_2
.\end{equation*}
Here, $N / A$ shifts the whole stress distribution uniformly up or down.
It is pure tension or compression.

The term $- M_2 / I_2 x_1$, creates a linear stress variation in $x_1$ caused by bending in about $e_2$ and similarly for $M_1 / I_1 x_2$.
